\documentclass[10pt,a4paper]{amsart}
\usepackage[margin=19mm]{geometry}
\usepackage{amsmath,amssymb,mathtools}
\usepackage[scr=rsfs,cal=euler]{mathalfa}
\usepackage{xfrac}
\usepackage[unicode,hidelinks,bookmarks=false]{hyperref}
\newcommand{\ie}{i.e.\ }
\newcommand{\cf}{cf.\ }
\newcommand{\resp}{respectively\ }
\newcommand{\Subsection}{\subsection}
\providecommand{\nobreakdash}{\nobreak\hbox{-}}
\setlength{\emergencystretch}{2em}
\multlinegap0pt
\usepackage{bm}
\usepackage{bbm}
\usepackage{dsfont}
\usepackage{xspace}
\usepackage{cancel}
\usepackage{mathtools}
\usepackage{cleveref}
\usepackage{amsthm}
\makeatletter
\@ifundefined{lemma}{\newtheorem{lemma}{Lemma}}{}
\makeatother
\title[Numerical Periodic Normalization at Codim 1 Bifurcations of ]{Numerical Periodic Normalization at Codim 1 Bifurcations of Limit Cycles in DDEs}
\author{M. M. Bosschaert, B. Lentjes, L. Spek, Yu. A. Kuznetsov}
\date{Source: arXiv:2505.19786v1 (CC BY 4.0)}
\begin{document}
\maketitle

\noindent\emph{Source and licence.} Numerical Periodic Normalization at Codim 1 Bifurcations of Limit Cycles in DDEs. Version arXiv:2505.19786v1 (CC BY 4.0), distributed under the Creative Commons Attribution 4.0 International licence, as recorded in the arXiv licence metadata for this exact version and shown on the abstract page https://arxiv.org/abs/2505.19786v1. Full rights notice: licenses/arxiv-2505.19786-CC-BY-4.0.txt.\par\smallskip

\noindent\textit{Editorial scope.} Counted unit: the Lemma whose source label is \texttt{lemma:adjoint charac} in the article (source0.tex lines 297--310, source label \texttt{lemma:adjoint charac}), delivered together with the article's own complete original proof of it (source0.tex lines 311--328, one \texttt{proof} environment of 18 lines). No mathematics is written, repaired or re-derived: statement and proof are the article's own text line for line. Required context retained uncounted: lemma:Delta(z)closed (lines 212--220); eq: pairingTCn (lines 283--285); lemma:nondegenerate (lines 288--290). Every definition, display and label of those lines is retained unchanged. Omitted from this package and not redistributed: 1--211, 221--282, 286--287, 291--296, 329--1221. Nothing inside a retained excerpt is omitted or altered: every retained body is delivered exactly as the article's lines are written, so no omission receipt of any kind accompanies this package. Citations: the retained excerpts carry no citation command at all, so no bibliography entry of the article is transcribed and no reference is redistributed. Reference adaptation: none was needed and none was made; every reference inside the delivered unit resolves inside it, so no unresolved reference or citation is produced. Printed slips kept literal: no printed word, symbol, hypothesis or number of the counted statement or of its proof was altered. Layout and class: the article's own class is \texttt{myarticle} with packages including listings, amsfonts, graphicx, amsmath, amsopn, physics, mathtools, thmtools, enumitem; the wrapper is the lane's pinned \texttt{amsart} editorial wrapper at 10pt on A4, which restates the theorem-like environments the retained text needs and adds the article's own macro definitions that the retained text needs (none), with extra wrapper packages (bm, bbm, dsfont, xspace, cancel, mathtools, cleveref). The delivered numbering is the unit's own. Assets: no retained span contains a figure environment, a graphics-inclusion command, a PDF-page inclusion, a table or a TikZ picture; no image file is copied into this package, the package declares no asset dependency and no image file is redistributed with it. Licence: version arXiv:2505.19786v1 of this article is distributed under the Creative Commons Attribution 4.0 International licence, as recorded in the arXiv licence metadata for this exact version and shown on the abstract page https://arxiv.org/abs/2505.19786v1. A full rights notice is delivered at licenses/arxiv-2505.19786-CC-BY-4.0.txt. Characters: every retained excerpt is pure ASCII, so the delivered engine, XeLaTeX with its default Latin Modern text font, reports no missing character.
\begin{lemma} \label{lemma:Delta(z)closed}
The operator $\Delta(z)$ is closed and $\rho(\Delta(z)) \neq \emptyset$ for all $z \in \mathbb{C}$, and there holds
\begin{align} 
\begin{split} \label{eq:not inv is floquet}
    \{ \sigma \in \mathbb{C} : \Delta(\sigma) \mbox{ is not invertible} \} &= \{ \sigma \in \mathbb{C} : \sigma \mbox{ is a characteristic value of } \Delta \} \\
    & = \{\sigma \in \mathbb{C} : \sigma \mbox{ is a Floquet exponent} \}.    
\end{split}
\end{align}
\end{lemma}

\begin{equation} \label{eq: pairingTCn}
    \langle p,q \rangle_T \coloneqq \int_0^T  p(\tau)q(\tau) d \tau, \quad \forall (p,q) \in C_T(\mathbb{R},\mathbb{C}^{n \star}) \times C_T(\mathbb{R},\mathbb{C}^n),
\end{equation}

\begin{lemma} \label{lemma:nondegenerate}
The bilinear map $\langle \cdot, \cdot \rangle_T$ from \eqref{eq: pairingTCn} is nondegenerate.
\end{lemma}

\begin{lemma} \label{lemma:adjoint charac}
For any $z \in \mathbb{{C}}$, the operator $\Delta^\star(z) : \mathcal{D}(\Delta^\star(z)) \coloneqq C_T^1(\mathbb{R},\mathbb{C}^{n \star}) \subseteq C_T(\mathbb{R},\mathbb{C}^{n \star}) \to C_T(\mathbb{{R}},\mathbb{C}^{n \star})$ defined by
\begin{equation} \label{eq:adjoint charac operator}
    (\Delta^\star(z)p)(\tau) \coloneqq -\dot{p}(\tau) + zp(\tau) - \int_0^h   p(\tau + \theta) e^{-z \theta} d_2 \zeta(\tau + \theta,\theta), \quad \forall p \in C_T^1(\mathbb{R},\mathbb{C}^{n \star}), 
\end{equation}
is the unique linear operator satisfying
\begin{equation} \label{eq:pairing charac}
    \langle \Delta^\star(z)p,q \rangle_T = \langle p,\Delta(z)q \rangle_T, \quad \forall (p,q) \in C_T^1(\mathbb{R},\mathbb{C}^{n \star}) \times C_T^1(\mathbb{R},\mathbb{C}^n).
\end{equation}
Moreover, $\Delta^\star(z)$ is closed and $\rho(\Delta^\star(z)) \neq \emptyset$ for all $z \in \mathbb{C}$, and there holds
\begin{equation*} 
    \{ \sigma \in \mathbb{C} : \Delta^\star(\sigma) \mbox{ is not invertible} \} = \{ \sigma \in \mathbb{C} : \sigma \mbox{ is a characteristic value of } \Delta^\star \}. 
\end{equation*}
\end{lemma}

\begin{proof}
Let $z,p$ and $q$ be given as mentioned. It follows immediately that
\begin{equation*}
    \langle p, \Delta(z)q \rangle_T = \int_0^T  p(\tau) \dot{q}(\tau) d\tau +  \int_0^T  zp(\tau)q(\tau) d\tau - \int_0^T p(\tau) \int_0^h d_2 \zeta(\tau,\theta) e^{-z\theta}q(\tau-\theta) d\tau.
\end{equation*}
Applying integration by parts on the first term of the right-hand side yields $-\langle \dot{p},q \rangle_T$, due to the $T$-periodicity of $p$ and $q$. The second term equals $\langle zp,q \rangle_T$, while the third term requires a bit more work. Let $P_m = \{0 = \theta_0 < \theta_1 < \dots < \theta_m = h \}$ be a partition of $[0,h]$ and choose an arbitrary $c_i \in [\theta_{i},\theta_{i+1}]$ for all $i = 0,\dots,m-1$. Consider  $\int_0^T  p(\tau) \sum_{i=0}^{m-1} [ \zeta(\tau,\theta_{i+1}) - \zeta(\tau,\theta_i)] e^{-z c_i}q(\tau - c_i)  d\tau$. If we substitute $ \tau - c_i \to \tau$, the integral bounds become $-c_i$ and $T - c_i$. Because $p,q$ and $\tau \mapsto \zeta(\tau,\cdot)$ are $T$-periodic, the integrand is $T$-periodic, and so we can shift the bounds of integration back to $0$ and $T$. Hence,
\begin{align*}
&\int_0^T p(\tau) \sum_{i=0}^{m-1} [ \zeta(\tau,\theta_{i+1}) - \zeta(\tau,\theta_i)] e^{-z c_i}q(\tau - c_i)  d\tau \\
&= \int_0^T \sum_{i=0}^{m-1} p(\tau+c_i) e^{-z c_i}[ \zeta(\tau + c_i,\theta_{i+1}) - \zeta(\tau + c_i,\theta_i)] q(\tau) d\tau.  
\end{align*}
If we let $m \to \infty$ while $\max_{i=0,\dots,m-1} |\theta_{i+1} - \theta_i| \to 0$, then the summation on the left-hand side approaches $\int_0^h d_2 \zeta(\tau,\theta) e^{-z\theta}q(\tau-\theta)$ by definition of the Riemann-Stieltjes integral. Hence,
\begin{equation*}
\int_0^T p(\tau) \int_0^h d_2 \zeta(\tau,\theta) e^{-z\theta}q(\tau-\theta) d\tau = \int_0^T \int_0^h p(\tau + \theta) e^{-z\theta}  d_2\zeta(\tau+\theta,\theta)q(\tau) d\tau,
\end{equation*}
and so $\langle p,\Delta(z)q \rangle_T = \langle \Delta^\star(z)p, q \rangle_T$, which shows that $\Delta^\star(z)$ from \eqref{eq:adjoint charac operator} satisfies \eqref{eq:pairing charac}. 

It remains to show that $\Delta^{\star}(z)$ is the unique linear operator satisfying \eqref{eq:pairing charac}. Let $Q(z) : C_T^1(\mathbb{R},\mathbb{C}^{n \star}) \to C_T(\mathbb{R},\mathbb{C}^{n \star})$ be a linear operator satisfying $\langle (Q(z) - \Delta^\star(z))p,q \rangle_T = 0$, for all $z \in \mathbb{C}$ and $ (p,q) \in C_T^1(\mathbb{R},\mathbb{C}^{n \star}) \times  C_T^1(\mathbb{R},\mathbb{C}^n)$. Let $q \in C_T(\mathbb{R},\mathbb{C}^n)$ be given and let $(q_m)_m$ be a sequence in $C_T^1(\mathbb{R},\mathbb{C}^n)$ such that $q_m \to q$ in norm as $m \to \infty$. As $\langle (Q(z) - \Delta^\star(z))p,q_m \rangle_T = 0$ for all $m \in \mathbb{N}$, we obtain $|\langle (Q(z) - \Delta^\star(z))p,q \rangle_T| \leq T \| (Q(z) - \Delta^\star(z))p \|_\infty \| q_m - q \|_\infty \to 0$ as $m \to \infty$ and thus $\langle (Q(z) - \Delta^\star(z))p,q \rangle_T = 0$ for all $q \in C_T(\mathbb{R},\mathbb{C}^n)$. But due to \cref{lemma:nondegenerate}, the map $\langle \cdot , \cdot \rangle_T$ is nondegenerate and so $Q(z) = \Delta^\star(z)$ on $C_T^1(\mathbb{R},\mathbb{C}^{n \star})$, i.e. $\Delta^\star(z)$ is uniquely determined. The remaining properties of $\Delta^\star(z)$ can be proven as performed in \cref{lemma:Delta(z)closed}.
\end{proof}

\end{document}
